% Français 9117, batch 3 — Gallica views 41–60.
% Pass: Opus 5 (claude-opus-5), 2026-09-14 — first pass, unchecked against
% the views by a human.
%
% What the twenty views hold. Twenty written pages, no blank and no plate:
% one continuous stretch of the prose essay the catalogue calls
% « Considérations générales sur l'état des sciences et des lettres », in a
% single hand throughout, a copy cleaner than a draft but corrected as it
% goes — words struck and rewritten in the line or above it, a few long
% passages cancelled, square brackets drawn in the margin before several
% paragraphs. No mathematics set as formulae. View 41 opens in the middle
% of a sentence (« logiques ? ») carried over from view 40; view 60 breaks
% off in the middle of one (« la cause d'une erreur »).
%
%   v. 41–42   Whether the observer's intellectual forms can bend their
%              objects: the history of the sciences shows a thousand errors;
%              each system rested on one certain truth, the rest filled in by
%              supposition; the inventors knew the difference, the schools
%              and the commentators did not.
%   v. 43–44   How, if logical necessities are absolute, a sound mind
%              reasoning from a certain principle can reach false
%              conclusions: the fault lies with language, made for usual
%              intercourse, and with technical terms that each school read
%              its own way.
%   v. 45–48   The history: the mathematical sciences, from their birth,
%              realised the « type du vrai »; Descartes doubts the schools
%              and rebuilds the world on the vortices, gives reason its
%              independence, joins algebra to geometry; Newton makes the
%              order of the universe a mathematical truth; the contrast
%              between the precise language of calculation and the
%              mysterious language of philosophy.
%   v. 49–51   Bacon; the ancients observed little; observation takes the
%              place of the « pourquoi » with the « comment » and then the
%              « combien »; physics becomes a branch of the exact sciences,
%              and the elements of algebra enter education.
%   v. 51–54   Back to the two questions: why so many errors (the type of
%              the true is abstract and cannot give particular realities;
%              the probability of guesses); the absolutism of logical
%              necessities vindicated by the mathematical theories, whose
%              language forbids wandering and betrays error; nature
%              confirming what analysis derives.
%   v. 55–58   « Les préliminaires qu'on vient de lire » — the turn to the
%              essay's own thesis: d'Alembert's sentence on the universe as
%              a single fact, quoted « déjà citée »; the fact is necessary;
%              essence and necessity are the same; contingent events and
%              the calculus of probabilities as a measure of ignorance
%              (a machine stopping at a given instant); the distinction
%              between contingent and necessary is the one between facts
%              whose nature is unknown and known; God's free will and the
%              balance whose heavier weight decides.
%   v. 59–60   Deliberation gives the feeling of liberty; necessity governs
%              the world; the unity of being, the dream of antiquity;
%              scepticism and « je pense, donc je suis »; innate ideas and
%              sensations each true within limits; man « l'abrégé de
%              l'univers ».
%
% Cancelled passages, transcribed inside \struck{} with their own
% corrections laid flat and described in a \note{}: a paragraph at view 53
% (on the absolutism of logical necessities, rewritten in the next
% paragraph), and at the head of view 54 the end of a sentence begun at
% the foot of view 53 (« et si quelque [fois] ») about an author who
% appropriates another's inventions — cancelled by a cross.
%
% The numbering on the leaves. Two series, both set out once in the
% opening \note{}. (1) Top left of every page, one number per page, struck
% through, in her ink: 31 to 50 across views 41 to 60, read as 31, 37, 41,
% 43, 46 and 50 at views 41, 47, 51, 53, 56 and 60 with some confidence,
% the others (32–36, 38–40, 42, 44–45, 47–49) cramped under the stroke and
% read in sequence rather than one by one; whether those at views 42 and
% 44 are struck could not be settled. (2) Top right of the rectos only, in
% a thin grey stroke unlike her black ink, the Bibliothèque's foliation:
% 20, 21, 22, 23, 24, 25, 26, 27, 28, 29 at views 41, 43, 45, 47, 49, 51,
% 53, 55, 57, 59, written \folio{20r} … \folio{29r} after \page{} on each
% recto; the versos carry no \folio{}. Beside 21 (view 43) a struck « 8 »
% and beside 27 (view 55) a second, struck « 27 », both in her ink. Every
% figure was read on the view; none is inferred from the view numbers.
%
% Notation and conventions. Her underlining is set \emph{} (« pourquoi »,
% « comment », « combien », « préliminaires », « à priori »). The spelling
% of the page stays: « tems », « connoître », « -oit », « enfans »,
% « datter », « rejetter », « essort », « satisfesoit », « clairvoyance »;
% also her grammar where it slips (« Le systèmes … satisfesoit » v. 41,
% « ses progrès … se répand » v. 59, « nous délibérions » v. 58, « Aussi
% voyons que une personne » v. 59, « L'une et l'autre opinions » v. 60),
% noted once at the spot. Where she writes a correction above a struck
% word the struck word comes first, then the correction. The marginal
% square brackets are recorded in notes. The d'Alembert sentence, marked
% with guillemets at the head of each line, is set as a quote.
%
% The hand. A clear, regular, sloping hand, the same from view 41 to 60;
% corrections in a smaller hand between the lines. Nothing in it points to
% a date.
%
% What could not be read (every \ill{}): v. 41, a two-sign interlinear
% addition above a struck « Si »; v. 44, the struck word after « La
% difficulté dont nous » and a short struck word before « Les »; v. 45, a
% flourish-word after a paragraph; v. 47, a short word above struck
% « fut »; v. 49, a flourish-word after « comparaisons », and a first
% struck correction above « l'envie d'imposer »; v. 50, a blotted sign in
% the left margin; v. 51, the struck word after « confus »; v. 52, the
% struck word in « l'analyse de … la théorie des probabilités »; v. 57, the
% two struck words after « fraction »; v. 58, the word under « rencontrent »,
% smeared out; v. 60, the pencil line at the head (only « Ch 4 » and
% « type » read), the struck word under « vraies », and a struck word
% before « dans tous les objets ». Offered with doubt (\uncertain{}): « on »
% v. 42; « l' » v. 44; « L'astronomie » v. 48; « diroit », « à » v. 49;
% « œil » v. 55 (under a blot); « Deracke » v. 55, a large underlined word
% after the sentence, not understood; « même » v. 53; « Tâchons » v. 59.
%
% Nothing here was checked against any published transcription.
\documentclass[11pt,a4paper]{article}
% The preamble shared by every transcription of the mathematicians' manuscripts in Gallica.
%
% It rests on one idea: **uncertainty compounds, it does not resolve.** A
% transcription that smooths over a doubtful word has destroyed the only thing
% it was made for — a reader must be able to tell, at every point, what was on
% the page and what was supplied. The apparatus macros below are therefore the
% critical apparatus, not decoration.
%
% The same commands are recognised by `scripts/render.mjs`, which produces the
% site's reading view, and by `scripts/tei.mjs`. A macro added here must be
% added there too, or rendering fails loudly — which is the wanted behaviour.
%
% Comments here are English, like the rest of the repository. The documents
% this preamble sets are French, and that is deliberate: see the skills.

\ProvidesPackage{germain}[2026/09/15 Transcriptions of mathematicians' manuscripts in Gallica]

\usepackage{amsmath,amssymb,amsthm}
\usepackage{mathrsfs}
% Kept from the parent preamble so the renderer's diagram support stays
% available; these manuscripts draw no commutative diagrams, and nothing here uses it.
\usepackage{tikz-cd}
\usepackage[english,french]{babel}
\usepackage{fontspec}

% --- Greek ------------------------------------------------------------------
% The text face stays fontspec's default, Latin Modern, which has no Greek:
% XeTeX drops every glyph it lacks with a line in the log and nothing on the
% page, and the Greek words the manuscripts quote vanished from the PDFs. So
% the Greek and Greek Extended blocks get a XeTeX character class of their own,
% and entering or leaving a run of it switches the family to CMU Serif — the
% same Computer Modern design, with polytonic Greek — and back. Only the
% family changes: a Greek word in \textit or \textbf keeps its shape and series.
% The transitions to and from every other class carry the tokens that class
% has towards an ordinary letter, so French spacing before « ; » or inside
% « » still holds after a Greek word. No grouping, so no brace or \endgroup
% can come out unbalanced; maths is left alone, since XeTeX inserts nothing
% there. Kept identical in `exercices.sty`.
\makeatletter
\newfontfamily\germainGreekfont{cmun}[
  Extension=.otf, NFSSFamily=germaingreek,
  UprightFont=*rm, ItalicFont=*ti, SlantedFont=*sl,
  BoldFont=*bx, BoldItalicFont=*bi, BoldSlantedFont=*bl]
\newXeTeXintercharclass\germain@greek
\def\germain@greekrange#1#2{%
  \@tempcnta=#1\relax
  \loop
    \XeTeXcharclass\@tempcnta=\germain@greek
    \ifnum\@tempcnta<#2\advance\@tempcnta\@ne\repeat}
\germain@greekrange{"0370}{"03FF}
\germain@greekrange{"1F00}{"1FFF}
\def\germain@greekprev{\rmdefault}
\def\germain@greekin{%
  \edef\germain@greekprev{\f@family}\fontfamily{germaingreek}\selectfont}
\def\germain@greekout{\fontfamily{\germain@greekprev}\selectfont}
\def\germain@greekjoin{%
  \XeTeXinterchartokenstate=\@ne
  \@tempcnta=\z@
  \loop
    \ifnum\@tempcnta=\germain@greek\else
      \XeTeXinterchartoks\germain@greek\@tempcnta=\expandafter{%
        \expandafter\germain@greekout\the\XeTeXinterchartoks\z@\@tempcnta}%
      \XeTeXinterchartoks\@tempcnta\germain@greek=\expandafter{%
        \the\XeTeXinterchartoks\@tempcnta\z@\germain@greekin}%
    \fi
    \ifnum\@tempcnta<4095\advance\@tempcnta\@ne\repeat}
% babel-french sets its own classes, and saves and restores the interchar
% state, each time a language is selected: join again after it does.
\addto\extrasfrench{\germain@greekjoin}
\addto\extrasenglish{\germain@greekjoin}
\AtBeginDocument{\germain@greekjoin}
\makeatother

\usepackage{geometry}
\geometry{a4paper,margin=25mm,marginparwidth=28mm}
\usepackage{marginnote}
\usepackage{xcolor}
\usepackage[normalem]{ulem}
\setlength{\emergencystretch}{3em}
\usepackage{hyperref}
\hypersetup{colorlinks=true,linkcolor=black,urlcolor=black,pdfborder={0 0 0}}

% --- Batch metadata ---------------------------------------------------------
% Taken as they stand by the rendering script, which shows them at the head of
% the reading view. They fix what the file claims to cover. The macro names are
% the parent project's, so that the scripts stay shared; \folder here means the
% volume's slug — `fr-9115` — and \pages counts Gallica views.

\newcommand{\folder}[1]{\gdef\grothFolder{#1}}
\newcommand{\batch}[1]{\gdef\grothBatch{#1}}
\newcommand{\pages}[2]{\gdef\grothFirst{#1}\gdef\grothLast{#2}}
\newcommand{\foldertitle}[1]{\gdef\grothFoldertitle{#1}}

% The holder's own shelfmark, printed as they write it: « Français 9115 ».
\newcommand{\shelfmark}[1]{\gdef\grothShelfmark{#1}}

% The dating, copied verbatim from the holder's catalogue — never inferred
% here. « XIXe siècle » is the BnF's; a pass that dates a leaf from its content
% says so in a \note{}, as its own inference.
\newcommand{\dating}[1]{\gdef\grothDating{#1}}

% The status notice, printed in the title block and repeated as a diagonal
% watermark on every page of the PDF. The manuscripts are in the public
% domain; what this line declares is not a right but a quality — an
% unverified first pass by a machine. The skills write the line; a file
% without it gets no watermark, so leaving it out is visible in the output.
\newcommand{\watermark}[1]{\gdef\grothWatermark{#1}}

\gdef\grothFolder{?}\gdef\grothBatch{}\gdef\grothFirst{?}\gdef\grothLast{?}%
\gdef\grothFoldertitle{}\gdef\grothShelfmark{}\gdef\grothDating{}\gdef\grothWatermark{}

% --- The résumé -------------------------------------------------------------
\newenvironment{resume}
  {\small\setlength{\parskip}{0.45em}}
  {\par\medskip\hrule\medskip}

% --- Keywords ---------------------------------------------------------------
% In English, closing the modernised reading's résumé: the modern vocabulary
% under which to search for what the leaves construct. The manifest extracts
% them to tag the volume — this line is the single source of the tags.
\newcommand{\keywords}[1]{\par\smallskip\noindent
  {\footnotesize\textsc{Keywords} --- \textit{#1}}\par}

% --- The critical apparatus -------------------------------------------------

% The Gallica view number, in the margin. This is the anchor that lets a
% disagreement over a formula be settled by looking at *one* image rather than
% a volume of 750. The site uses it to turn the facsimile as the transcription
% is scrolled. A view is one scanned image: usually one side of a leaf.
\newcommand{\page}[1]{%
  \marginnote{\footnotesize\color{gray!70}\textsf{v.\,#1}}%
  \label{page:#1}\ignorespaces}

% The BnF's pencilled foliation, where the leaf carries it and a pass can read
% it: « 198r ». This is what the literature cites — « ff. 198r–208v » — and
% the one line that turns a citation into a view. Recorded, never inferred:
% a folio a pass cannot read is not written.
\newcommand{\folio}[1]{%
  \marginnote{\footnotesize\color{gray!50}\textsf{f.\,#1}}\ignorespaces}

% The modernised reading's coarser counterpart to \page{}: which views a
% section is drawn from.
\newcommand{\pagerange}[2]{%
  \marginnote{\footnotesize\color{gray!70}\textsf{v.\,#1--#2}}%
  \ignorespaces}

% Illegible: never guessed.
\newcommand{\ill}{\textcolor{red!60!black}{[\,\ldots\,]}}

% A doubtful reading: the word is given, and flagged as uncertain.
\newcommand{\uncertain}[1]{\uline{#1}}

% An editorial addition — a missing letter, an implied word.
\newcommand{\add}[1]{\textcolor{blue!50!black}{[#1]}}

% Struck out by the writer. What was crossed out is part of the document.
\newcommand{\struck}[1]{\sout{#1}}

% The transcriber's note, on the state of the page or the mathematics.
\newcommand{\note}[1]{\par\smallskip{\small\color{gray!60!black}\textit{[#1]}}\par\smallskip}

% A marginal note of hers, distinct from the body of the text.
\newcommand{\marginal}[1]{\par\smallskip\noindent\hspace{1em}%
  {\small\color{gray!60!black}#1}\par\smallskip}

% --- Title ------------------------------------------------------------------
% The title block is French, like the documents themselves.

\AddToHook{shipout/background}{%
  \ifx\grothWatermark\empty\else
    \begin{tikzpicture}[remember picture, overlay]
      \node[rotate=52, scale=3.4, align=center, text=gray!18,
            font=\sffamily\bfseries]
        at (current page.center) {\grothWatermark};
    \end{tikzpicture}%
  \fi}

\AtBeginDocument{%
  \begin{center}
    {\large\bfseries Manuscrits de mathématiciens (Gallica) ---
      \ifx\grothShelfmark\empty\grothFolder\else\grothShelfmark\fi}\\[2pt]
    {\itshape\grothFoldertitle}\\[4pt]
    {\small vues \grothFirst--\grothLast
      \ifx\grothBatch\empty\else\ (lot \grothBatch)\fi}\\[2pt]
    \ifx\grothDating\empty\else
      {\small\color{gray!60!black}Datation du catalogue : \grothDating}\\[2pt]
    \fi
    \ifx\grothWatermark\empty\else
      {\small\bfseries\color{red!45!gray}\grothWatermark}\\[2pt]
    \fi
    {\footnotesize\color{gray!60!black}Transcription établie d'après les images IIIF de Gallica.
     Source gallica.bnf.fr / Bibliothèque nationale de France.\\
     Les lectures incertaines sont soulignées, les illisibles notées [\,\ldots\,],
     les ajouts entre crochets ; \textsf{v.} numérote les vues de Gallica,
     \textsf{f.} la foliotation au crayon de la BnF.}
  \end{center}
  \vspace{1em}}

\endinput


\folder{fr-9117}
\shelfmark{Français 9117}
\batch{3}
\pages{41}{60}
\dating{1801-1900}
\watermark{Lecture automatique\\non vérifiée}
\foldertitle{« Considérations générales sur l'état des sciences et des lettres aux différentes époques de leur culture, » par Sophie GERMAIN.}

\begin{document}

\note{Numérotation des pages. Deux séries. En haut à gauche de chaque
page, un nombre biffé, à l'encre, de 31 (vue 41) à 50 (vue 60), un par page ; lus
avec assurance : 31, 37, 41, 43, 46, 50 aux vues 41, 47, 51, 53, 56, 60 ;
les autres sont serrés sous le trait et lus dans la suite. En haut à
droite des rectos seulement — vues 41, 43, 45, 47, 49, 51, 53, 55, 57,
59 — les nombres 20 à 29 ;
à côté de 21 (vue 43) un « 8 » biffé, à côté de 27 (vue 55) un second
« 27 » biffé, tous deux à l'encre de l'auteur. Cette seconde série, d'un
trait fin et gris, est la foliotation de la Bibliothèque ; elle est portée
en « f. » à chaque recto. La vue 41 commence au milieu d'une
phrase commencée à la vue 40 ; la vue 60 s'arrête au milieu d'une
phrase.}

\page{41}\folio{20r}
logiques ? Les formes intellectuelles de l'observateur auroient-elles
donc le pouvoir de ployer à leurs convenances les sujets soumis à son
examen ?

Nous sommes loin de cette heureuse position. L'histoire des sciences
nous présente mille écarts, et l'esprit humain a employé plus d'efforts
à \struck{renverser} détruire ses propres ouvrages qu'à en reconstruire
de nouveaux. \struck{Chaque} Le systèmes satisfesoit, à l'aide des
suppositions les plus hasardées, aux faits qu'il\struck{s}
devoit\struck{ent} expliquer ; \struck{bientôt} \struck{ils} ces systèmes
devenoient insuffisans, mais leur influence sur l'esprit des philosophes
étoit alors un obstacle \struck{qu} difficile à vaincre pour arriver à
la connoissance de la vérité.\note{« Le » est écrit au-dessus de
« Chaque » biffé, mais « systèmes » garde son \emph{s} ; le verbe et
« il\struck{s} devoit\struck{ent} » sont au singulier. « ces systèmes »
est ajouté au-dessus de « ils » biffé ; « mais » est serré dans
l'interligne.}

\struck{Si} \ill{}. Si\note{Au-dessus de « Si » biffé, une addition de
deux signes, soulignée, dont le premier n'a pas été lu ; le second est
« Si ».} nous portons en nous-mêmes le type du vrai, pourquoi avons-nous
commis tant de méprises ! Si notre logique n'est autre chose que le
recueil des préceptes de la raison absolue, comment, \struck{avec}
malgré les secours d'un guide sûr, avons nous pû errer si longtems dans
la région nébuleuse des suppositions gratuites.

En examinant la première de ces questions, Nous verrons qu'il est hors
de doute que le type du vrai n'a jamais cessé de se faire sentir au
milieu des écarts de la raison. Et, en effet, chaque système a été
inspiré par la connaissance d'une vérité incontestable. L'homme de
génie, frappé de l'importance de cette vérité, et persuadé de l'unité
de l'être, a voulu rapporter toutes choses à celle dont il avoit acquis
la certitude. Alors il a imaginé, il a supposé, il a rempli, d'une
manière plus ou moins heureuse,

\page{42}
les nombreux intervalles entre les points solidement établis. Mais cet
esprit supérieur auteur d'un vaste système n'a jamais confondu, dans sa
conscience, la certitude absolue qui servoit de base à l'œuvre de la
pensée avec la probabilité, souvent bien foible à ses propres yeux, des
suppositions destinées à lier entre elles les diverses parties de sa
doctrine.

C'est là c'est dans la pensée des inventeurs qu'il faut étudier la
nature de l'intelligence humaine, et,\note{« et, » est ajouté dans la
marge de gauche.} à cet égard l'histoire du genre humain se réduit à
celle d'un petit nombre d'hommes nés avec l'honorable mission d'éclairer
leurs semblables.

Sans doute le modèle du vrai, qui nous sert à reconnoître le bon et le
beau, n'est pas le partage exclusif de ces hommes privilégiés ; mais des
esprits communs ne voyent qu'autour d'eux. Dans le cercle de leurs
affections et de leurs intérêts, ils sont juges éclairés ; au delà il
n'existe aucune certitude dont ils fassent cas ; et d'ailleurs ils
manqueroient de facultés pour savoir l'apprécier.

Revêtues des formes séduisantes qu'une imagination élevée sait prêter à
ses conceptions, les doctrines systématiques ont été adoptées avec
enthousiasme par la curiosité publique ; les esprits cultivés en ont
fait leur pâture ; elles ont été enseignées dans les écoles. Il
s'agissoit de les savoir, et non pas de les juger. On en
suivoit les conséquences ; on en multiplioit les applications ; tout le
monde parloit d'après le maître, on expliquoit sa pensée ;
\uncertain{on} faisoit sur ses écrits mille commentaires qu'\struck{il}
lui même n'eût \struck{sans} certainement pas adoptés.\note{« lui même »
est ajouté au-dessus de la ligne, avec un signe d'insertion, et
« certainement » au-dessus de « sans » biffé. Le mot lu « on », après
« pensée ; », est souligné et surchargé.}

\page{43}\folio{21r}
La simplicité primitive disparoissoi\struck{en}t ; une foule d'erreurs
venoient obscurcir le fond de vérité qui avoit éclairé le premier
auteur du système ; et pourtant l'enseignement et le crédit y étoient
obstinément attachés, jusqu'à ce qu'une hypothèse plus en harmonie avec
les progrès de l'observation eût satisfait au besoin de savoir, qui a
devancé, pendant un tems si long, la création de la science véritable.

Ici se présente naturellement la question relative aux certitudes
logiques. Comment, si elles sont absolues, l'esprit humain a-t-il pû,
s'abandonner à l'erreur ?

Il est d'abord évident que tout faux raisonnement, dès lors qu'il peut
être jugé tel par la raison humaine, doit être attribué à toute autre
cause qu'au défaut d'absolutisme dans nos nécessités intellectuelles. Il
ne reste donc autre chose à examiner que les déviations commises, par
l'homme de bonne foi et de jugement éclairé, qui, partant d'un principe
certain et raisonnant avec la plus \struck{s} sévère exactitude, est
cependant arrivé à des conclusions démenties par les faits.

Nous observerons qu'il est extrêmement difficile d'exprimer le principe
certain dont on veut suivre les conséquences, d'une manière assez
précise pour être sûr qu'il soit compris tout entier dans la définition
à laquelle il donne lieu, et qu'en même tems, aucune idée qui

\page{44}
ne seroit pas nécessairement comprise dans le principe certain n'ait pas
été indûment introduite dans la définition.

La difficulté \struck{dont nous \ill{}} dont nous parlons tient à la
nature des langues. Elles doivent leur origine à des communications
usuelles ; dans les idées qui se rapportent aux choses présentes ou à
celles qui sont parfaitement connues, elles ont toute l'exactitude
désirable ; mais, pour les \struck{faisant servir au} rendre
applicables à des sujets philosophiques, il a fallu prendre au figuré
des termes fort clairs, à la vérité, dans leur signification propre,
mais dont la précision ne pouvoit se conserver lorsqu'on altéroit le
sens dans lequel on étoit accoutumé à les entendre.

Cette difficulté a toujours été sentie. On a crû l'éluder en forgeant
des mots nouveaux, lorsqu'on avoit à exprimer des idées nouvelles ; Il
est évident cependant que ces mots eux-mêmes \struck{étoient} avoient
besoin d'être définis, et ne remédioient nullement à l'inconvénient du
défaut de précision. Loin de là ; \struck{même}, les expressions
techniques ont été interprétées de manières diverses, suivant les
opinions des \struck{savans} personnes qui cherchoient dans un système
accrédité \struck{un} \uncertain{l'}appui \struck{en faveur} pour de
leurs idées particulières. \struck{Elles sont} les expressions
techniques sont ainsi devenues une des sources les plus fécondes, de la
divagation des opinions philosophiques. \struck{\ill{}} Les
\struck{mots} expressions étoient les mêmes ; \struck{mais} chacun avoit
une opinion disparate à celle de son interlocuteur.\note{« pour » est
écrit au-dessus de « en faveur » biffé, sans que « de » paraisse biffé ;
« les expressions techniques sont » est écrit au-dessus de « Elles
sont » biffé. Un mot court, biffé, se lit au-dessus de la ligne avant
« Les ».}

\page{45}\folio{22r}
Dans un tems reculé, dont il est sans doute difficile d'assigner les
premières époques \struck{commencement, l'esprit humain}, les propriétés
générales des nombres, celles des figures simples et des corps
réguliers, avoient attirées l'attention des hommes nés avec le génie des
sciences exactes. Ici les idées sont d'une extrême simplicité. On avoit
sous les yeux les figures et les corps géométriques eux mêmes ; il étoit
\struck{donc} impossible de leur attribuer des propriétés qu'ils
n'avoient pas.\note{« époques » est écrit au-dessus de la ligne, le
\emph{s} détaché ; « les » surcharge « le ».}

Des remarques multipliées ont conduit à la connoissance parfaite de ces
objets ; un grand nombre de théorèmes curieux en ont été le fruit,
\struck{et} lorsqu'on a voulu exprimer ces théorèmes et ceux qui
concernent les nombres, quelques signes dont la signification ne
pouvait être équivoque ont suffi pour représenter avec précision des
idées d'une exactitude parfaite.

Dès leur naissance, les sciences mathématiques ont offert à l'esprit
humain l'entière réalisation de ce type du vrai, objet de ses plus
chères affections. Partout ailleurs, il en cherchoit en vain les
caractères sublimes. Mille suppositions gratuites avoient été
incorporées à un petit nombre de vérités ; et, malgré les formes
absolues de l'enseignement philosophique, l'homme doué d'un esprit
juste, sentoit au fond de sa conscience que l'étude ne pouvoit le
conduire à aucune certitude véritable.\note{Un crochet ouvrant est tracé
devant « Partout ». Sous la fin du paragraphe, à droite, un mot en
paraphe, \ill{}.}

Les tems ne sont pas encore fort éloignés où les sciences physiques,
morales, religieuses, et politiques étoient \struck{obscurcies}
surchargées \struck{par} d'une foule de doctrines hypothétiques et
mystérieuses. Aussi voyons-nous qu'alors la géométrie inspiroit un
enthousiasme que nous ne retrouvons plus au même degré.

\page{46}
Et comment en effet, après s'être astreint à étudier avec application
les divers systèmes qui composoient la science, systèmes rendus plus
obscurs encore, par une foule de commentaires, dont les auteurs étoient
loin de la sagacité des premiers inventeurs et se contredisoient entre
eux de \struck{mille} cent manières diverses, l'homme doué du sentiment
profond des conditions qui n'appartiennent qu'au vrai, n'auroit-il pas
été transporté d'une joie indicible en trouvant au plus haut degré cette
\struck{évidence} clair\struck{e}voyance de la vérité, dont la privation
l'avoit si cruellement tourmenté ?

Descartes osa douter publiquement des doctrines de l'école.
\struck{mais} Le grand homme n'eut pourtant pas le courage de renoncer à
l'espérance \struck{mille} tant de fois deçue, d'arriver enfin à
réaliser la copie fidèle du type de l'être \struck{; ce type étoit
empreint avec force dans son génie créateur}, il reconstruisit l'univers
\struck{tout} sur un nouveau plan. Mais le noble exemple qu'il avoit
donné servit bientôt à faire rejetter son propre système.
Descartes\note{« Descartes » est écrit sur un « Il ».} rendit ainsi à la
raison \struck{publique} un service immense ; il créa pour elle une
époque nouvelle ; elle lui doit son indépendance. L'hypothèse ingénieuse
des tourbillons, semblant appartenir au tems qui venoit de finir, aussi
en marqua-t-elle la dernière limite, et les efforts de l'esprit humain
changèrent-ils alors entièrement de direction.

\page{47}\folio{23r}
Les sciences mathématiques étoient composées de deux parties
distinctes : Descartes sut les réunir. Elles avoient déjà fait d'assez
grands progrès : l'application de l'algèbre à la géométrie leur donna un
nouvel essort ; Elles étoient isolées de toutes autres recherches ; la
langue des calculs, déjà ployée à un usage nouveau, \struck{devint}
\struck{fut} bientôt après susceptible d'exprimer les grands faits du
ciel.\note{Au-dessus de « fut » biffé, un mot court, \ill{}, qui paraît
biffé lui aussi ; « après » est ajouté au-dessus de la ligne.} Ainsi le
même homme qui avoit \struck{déjà} eu la gloire de renverser d'anciennes
erreurs eut la gloire, plus grande encore, d'ouvrir à ses successeurs
une route dans laquelle il étoit impossible \struck{qu'ils
s'égarassent} de s'égarer.

Newton parut, armé d'un nouveau genre de calcul : \struck{et} l'unité,
l'ordre \struck{et} les proportions de l'univers, que le sentiment du
vrai avoit fait chercher si longtems, devinrent des vérités
mathématiques. Son génie avoit reconnu la cause des mouvemens célestes ;
une analyse pleine de finesse lui servit à les mesurer. L'optique devint
aussi entre ses mains une science nouvelle, et il devina, par rapport à
la nature des corps réfringens des vérités qu'il étoit réservé à la
chimie de vérifier, longtems après.

C'est de cette époque, à jamais mémorable, qu'il faut datter l'alliance
entre les sciences mathématiques et les sciences physiques. La mécanique
et l'hydrodynamique avoient été connues des anciens. De grands travaux
et les livres d'Archimède attestent qu'ils en savoient et la pratique et
la théorie. mais l'idée des quantités est tellement inhérente à celle
des forces qu'on peut dire de ces deux sciences qu'elles sont
essentiellement mathématiques.\note{Un crochet ouvrant est tracé devant
« La mécanique ».}

\page{48}
Leurs élémens, ceux de l'algèbre, et ceux de la géométrie composoient
tout le domaine des idées exactes. Partout ailleurs, on ne retrouvoit
plus que les vains efforts du génie pour arriver à la connoissance de la
vérité, et les erreurs sans nombre que les doctrines insuffisantes des
premiers inventeurs traînoient à leur suite. Le langage mystérieux
employé par les philosophes, langage plus obscur encore que les idées
qu'il étoit destiné à exprimer, formoit avec la langue précise et claire
des sciences exactes un contraste singulier.

Dans un tems où les géomètres vivoient isolés et où ils étoient en petit
nombre, ce contraste étoit connu d'eux seuls ; et son effet se bornoit à
leur inspirer le plus profond mépris pour toutes les autres sciences.
Mais lorsque les phénomènes célestes, objets d'admiration et de la
curiosité des hommes, vinrent se ranger sous les lois du calcul, l'étude
des mathématiques devint plus générale, \struck{et le contraste entre la
rigueur de leur méthode et sur l'argumentation stérile des école,
frappa tous} et les bons esprits furent frappés d'une manière
d'argumenter si différente de celle de l'école.\note{Les trois lignes
biffées portent en interligne « supériorité », biffé avec elles. « d'argumenter
si différente de celle de l'école » est écrit sous la ligne, en
interligne serré.}

\uncertain{L'astronomie} physique remplaçoit des hypothèses discréditées ;
une vive lumière succédoit à l'assemblage des idées les plus
\struck{bizarres}, obscures. Cette révolution subite ébranla l'empire des
préjugés ; elle alarma les hommes intéressés à en soutenir le règne. Ils
craignoient les vérités les plus étrangères à leurs doctrines ;
\struck{et} aucune profession de foi ne parut assez orthodoxe pour les
rassurer.\note{Un crochet ouvrant est tracé devant ce paragraphe.}

\page{49}\folio{24r}
semblables au peintre qui éloigne des regards du spectateur tout objet
réel et palpable, l'instinct d'une sorte de perspective morale les avoit
avertis du danger des comparaisons.\note{À la suite, au milieu de la
ligne, un mot en paraphe, souligné, \ill{}.}

Tandis que le système du monde présentoit aux philosophes, le spectacle
nouveau d'un mécanisme simple dans son principe et fécond dans ses
conséquences, la physique sublunaire étoit encore surchargée de mille
suppositions, nées du besoin d'expliquer les faits dont la liaison étoit
inconnue. L'attachement aux vieilles routines, lorsqu'il étoit exempt de
\struck{l'envie d'imposer} l'envie d'imposer, ne pouvoit tenir longtems
\struck{ne pouvoit tenir longtems} \struck{contre le désir} contre le
désir et l'espérance d'obtenir dans d'autres genres d'études des succès,
dont un grand exemple venoit de révéler la possibilité.\note{Au-dessus
de « l'envie d'imposer » biffé, une première correction, biffée elle
aussi, \ill{} ; « l'envie d'imposer » est récrit dans l'interligne, à
gauche, au-dessus de « ne pouvoit ».} Les sciences se composoient d'un
mélange confus d'erreurs et de vérités : on \struck{le} sentit qu'il
falloit tout refaire. Bacon en donna le conseil.

Les anciens guidés par des considérations métaphysiques avoient peu
observé, \struck{ils semblent avoir} on \uncertain{diroit} qu'ils ont
craint de rencontrer dans la réalité des faits le démenti \uncertain{à}
de leurs idées systématiques.\note{Un crochet ouvrant est tracé devant
« Les anciens ». « on \uncertain{diroit} qu'ils ont » est écrit dans
l'interligne, au-dessus de la partie biffée.} À la renaissance des
lettres on étudia leurs écrits, leur littérature offroit des modèles ;
elle obtint à juste titre l'admiration universelle. \struck{cette
admiration on l'accorda aussi à tout ce qu'ils avoient laissé} Tous
leurs ouvrages furent également recherchés, également admirés. On adopta
leurs idées ; et les controverses \struck{vagues} n'eurent d'autre objet
que les diverses manières dont il étoit possible de les interpréter. Si
quelquefois on essaya des explications nouvelles, ce fut toujours, à
leur exemple, en cherchant à ployer les faits à des explications
également vagues et hasardées.

\page{50}
Jusque là, on avoit toujours cherché les causes des phénomènes, on
commença alors à les considérer en eux mêmes. Au lieu du \emph{pourquoi}
on voulut savoir le \emph{comment} de chaque chose. Une foule
d'observateurs laborieux commencèrent à examiner la nature des faits.
Ils renoncèrent courageusement pour eux-mêmes à la satisfaction de les
expliquer ; dans l'espérance de léguer à leurs successeurs une masse de
\struck{faits} connoissances positives, dont la liaison se
\struck{présenteroient} montreroit nécessairement dans un tems plus
éloigné.

C'est de cette époque que datte la connoissance de la nature. Jusque là
l'homme l'avoit \struck{imaginée} ; il la vit alors pour la première
fois.\note{« alors » est encadré d'un trait. Dans la marge de gauche,
en face de cette ligne, un signe barré et taché d'encre, \ill{}.}

On chercha à mesurer tout ce qui est mesurable. À la question du
\emph{comment} se joignit celle du \emph{combien}. Les phénomènes mieux
appréciés devinrent calculables, les plus simples d'entre eux
\struck{devinrent entre les mains des} présentèrent aux successeurs de
Newton des objets de recherches. De nos jours, l'esprit mathématique a
fait de tels progrès que la physique dite particulière, c'est à dire la
connoissance des phénomènes naturels qui n'appartiennent pas à l'histoire
naturelle, a pour ainsi dire disparu, et s'est trouvée transformée en une
des branches les plus importantes des sciences exactes.

En se ployant aux usages nouveaux, la langue du calcul s'est enrichie de
plusieurs méthodes nouvelles ; et ces méthodes ont offert ensuite le
moyen de traiter des questions qui sembloient, il y a peu de tems, devoir
rester étrangères aux \struck{mathématiques} sciences exactes.

\page{51}\folio{25r}
De si grands progrès, des applications si nombreuses ont tourné tous les
esprits vers les sciences mathématiques. Il y a moins d'un siècle, leur
objet étoit circonscrit dans un petit nombre de vérités abstraites ; les
\struck{gens} personnes les plus instruites regardoient l'algèbre comme
un langage barbare, \struck{un géomètre on s'étonnoit} et
indéchiffrable. Aujourd'hui les élémens de cette science entrent dans
l'éducation, son esprit a pénétré dans la masse des nations, et
\struck{l'esprit} la raison publique y a puisé des forces
nouvelles.\note{« et indéchiffrable » est la fin de la ligne, la
dernière syllabe débordant dans la marge, avec un signe de renvoi.}

Nous venons de voir comment l'esprit humain, après s'être épuisé en
vains efforts pour réaliser au dehors de lui le modèle du vrai empreint
dans sa pensée, changeant tout à coup de direction, abandonna les
espaces vagues d'une \struck{metta}physique ténébreuse ; parcourut pas à
pas la route de l'observation ; et, profitant avec art des ressources
offertes par les progrès d'une science où s'étoient réfugiées les idées
d'ordre et de rectitude, qui partout ailleurs étoient ensevelies sous un
\struck{confus \ill{}} amas confus de théories bizarres et hétérogènes,
parvint à soumettre aux lois du calcul, des phénomènes \struck{qui} dont
la nature étoient restée longtems inconnue.\note{« dont la nature » est
écrit au-dessus de « qui » biffé ; « étoient » n'a pas été corrigé, mais
les finales de « restée » et « inconnue » sont récrites. « amas confus »
est écrit à la suite de la ligne et descend dans la marge.}

Reprenons présentement les deux questions que nous nous sommes
proposées. La digression historique à laquelle nous nous sommes livrée
nous fournira le moyen d'y répondre avec plus de précision.

On demande d'abord pourquoi nous avons commis tant de méprises, si nous
portons en nous mêmes le type du vrai.

\page{52}
Nous voyons clairement que l'esprit humain, pressé de jouir, avoit,
jusqu'à nos tems modernes, suivi une route où il ne devoit rencontrer
aucune réalité effective. Éclairés par l'expérience, il nous est même
facile aujourd'hui de voir \emph{à priori} pourquoi les faits
échappoient à chaque instant aux divers systèmes que le génie
\struck{avoit} de l'homme avoit enfantés.

Et, en effet, le type du vrai, par sa nature se compose d'idées
abstraites. Il nous avertit de ce qui répugne, c'est à dire de ce qui ne
peut exister en même tems ; mais il ne peut suffire pour nous manifester
des réalités particulières. Nous savons que chaque chose a son essence ;
cette essence est l'unité du sujet ; elle est susceptible de division ;
nous savons encore qu'il existe de l'ordre et des proportions entre ses
parties. Mais, dans un cas donné, quelle essence, quel ordre, quelles
proportions devons-nous rencontrer ? le modèle de l'être ne suffit pour
nous en informer.

Lorsqu'étayé par un petit nombre de connoissances certaines, l'homme de
génie a essayé de suppléer, \struck{par} en admettant des suppositions
gratuites, à ce qui lui manquoit d'observations positives, il a établi
entre ces choses des relations purement fantastiques : \struck{et si}
Dans l'état actuel de nos connoissances, on demandoit au géomètre
\struck{l'analyse de \ill{} la théorie des probabilités}, combien de
fois la théorie des probabilités \struck{veut} que des conjectures ainsi
formées\note{« géomètre » est écrit au-dessus de la partie biffée ; le
« D » de « Dans » surcharge une minuscule. Le mot lu « veut », en fin de
page, paraît biffé ; la phrase se poursuit à la vue suivante.}

\page{53}\folio{26r}
se \struck{trouvent en effet} réalisées, il rencontreroit bien
certainement la très petite fraction qui représente les succès obtenus
pendant les siècles qui ont précédé l'époque \struck{dont nous avons
rappelé la gloire}, où les observations précises ont remplacé les
assertions dénuées de preuves.\note{« se » surcharge un autre mot.}

\struck{À l'égard de la question relative à l'absolutisme de nos
nécessités logiques nous avons déjà répondu montré comment en partant
d'un principe certain il est possible cependant d'arriver à des
conséquences erronées, nous avons dit signalé trouvé des causes dans
l'ins inhabileté des langues formées pour des relations usuelles, à
représenter les idées et logiques abstraites avec une exactitude logique
parfaite.}\note{Ce paragraphe est annulé d'un long trait oblique. Il
porte en outre ses propres corrections, qu'on donne ici à plat :
« répondu », « montré », « comment », « logiques » (après
« nécessités »), « dit », « signalé », « l'ins », « et logiques » et
« logique » sont biffés ; « trouvé des causes d » est écrit au-dessus de
« signalé », et « abstraites » au-dessus de « et logiques ».}

À l'égard de l'objection contre l'absolutisme des nécessités logiques,
tirée d'exemples de démentis donnés par les faits à des conséquences
déduites d'un principe certain, nous avons dit comment l'imperfection des
langues introduisoit inopinément des idées étrangères au sujet ; en sorte
qu'on ne pouvoit être sûr ni de l'avoir fait entrer tout entier ni de lui
avoir adjoint aucun autre objet dans la définition qui sert de fondement
aux raisonnemens.

Cette explication est aujourd'hui pleinement justifiée par les théories
mathématiques, et l'absolutisme des nécessités logiques semble ne pouvoir
plus être revoqué en doute.

Par une suite d'efforts, concentrés cependant entre un bien petit nombre
d'hommes, une langue précise, exacte et où la moindre erreur deviendroit
sensible, a été formée et enrichie. Cette langue est celle de la raison
dans toute sa pureté. Elle \struck{repousse} interdit la divagation ;
elle signale l'erreur involontaire ; et il faudroit ne la pas connoître
pour essayer de la faire servir à l'imposture. \struck{et si
quelque}\note{Au-dessus de « quelque » biffé, un mot, lu
\uncertain{même}, qui paraît biffé aussi.}

\page{54}
\struck{fois il est arrivé qu'un auteur géomètre respectant cependant ait
voulu, dans une intention intérêt d'amour propre s'approprier des
inventions qui ne lui appartenoient pas il n'a pû même en copiant le
véritable auteur dissimuler longtems, que l'idée principale ne lui
appartenoit pas.}\note{Ces six lignes, qui achèvent la phrase laissée
en suspens au bas de la vue précédente, sont annulées d'une croix de deux
longs traits et soulignées d'un trait. Elles portent leurs propres
corrections : « auteur » est biffé sous « géomètre », « intention » sous
« intérêt », et « respectant cependant » est biffé.}

Elle reproduit dans toutes ses conséquences le principe qui lui a été
confié. Elle \struck{prouve} peut servir à prouver que l'unité
d'essence, l'ordre et les proportions du sujet, que l'esprit humain
cherche \struck{obstinément} dans tous les objets de son attention,
\struck{ne sont} n'expriment pas seulement \struck{un besoin dû à ses
formes intellectuelles} les conditions de notre satisfaction
intellectuelle, mais \struck{sont en effet celles de} qu'elles
appartiennent à l'être ou à la vérité.\note{Les corrections
interlinéaires de la fin de la phrase sont serrées : « n'expriment »
au-dessus de « ne sont », « qu'elles » au-dessus de « sont »,
« appartiennent », lui-même corrigé, au-dessus de « en effet », et
« à » au-dessus de « de ». « \emph{la} vérité » a été récrit « la
vérité » sur un « de » effacé.}

Et en effet lorsqu'on parvient à rendre une question mathématique c'est
à dire lorsqu'on a eu l'art d'en saisir l'essence d'une manière assez
simple pour que l'analyse puisse s'en emparer, la nature, docile à la
voix de l'homme, vient sanctionner les oracles de la science. Un fait
connu, bien apprécié, s'étoit présenté à sa pensée comme une conséquence
d'un ordre de choses encore inconnu ; il a sû définir cet ordre, et
bien-tôt l'expérience, abondante en circonstances nouvelles, met en
évidence et le génie qui a deviné son existence et l'excellence de la
méthode qu'il a sû \struck{mettre en œuvre} employer.

Douteroit-on que le type de l'être ait une réalité absolue, lorsqu'on
voit la\note{Un crochet ouvrant est tracé devant « Douteroit-on ». La
phrase se poursuit à la vue suivante.}

\page{55}\folio{27r}
langue des calculs faire jaillir d'une seule réalité dont elle s'est
emparée, toutes les réalités liées à la première par une essence
commune ? comment si de telles liaisons n'avoient en leur faveur que la
faculté de notre intelligence pour les concevoir, comment arriveroit-il
que l'observation des faits vînt, par une voie si différente, montrer en
dehors de la pensée de l'homme, l'édifice \struck{correspondant à celui
que} semblable à celui dont il trouve le modèle au dedans de
lui-même ?\note{Au-dessus de « lui », un petit « 2 » ; à la suite de la
ligne, en plus grands caractères et souligné, un mot lu
\uncertain{Deracke}, dont le sens n'a pas été compris.}

Les \emph{préliminaires} qu'on vient de lire nous ont paru nécessaires
pour bien entendre les idées que nous allons exposer, ils en fixeront le
sens et serviront peut être à leur faire pardonner ce qu'elles semblent
avoir de hardiesse et de nouveauté.\note{« avoir » est ajouté au-dessus
de la ligne.}

Dans l'état actuel de notre culture intellectuelle, nous avançons vers
la réalisation de ce qui fut un pressentiment chez les auteurs de tant
de systèmes prématurés.

Ils vouloient ramener toutes choses à une seule ; ils cherchoient
l'unité de l'être dont la nécessité s'est toujours fait sentir aux
esprits supérieurs. Cette pensée constante des hommes qui forment, à
travers les siècles, la chaîne des idées nécessaires du genre humain a
été clairement exprimée par d'Alembert lorsqu'il a écrit cette phrase
déjà citée :

\begin{quote}
« L'univers, pour qui sauroit l'embrasser d'un seul coup
d'\uncertain{œil} seroit un fait unique, une grande vérité. »
\end{quote}
\note{« d'un seul coup d'\uncertain{œil} » est ajouté au-dessus de la
ligne ; le dernier mot est couvert d'une tache d'encre. Chaque ligne de
la citation s'ouvre par un guillemet dans la marge.}

Ajoutons que, suivant notre conviction intime, ce fait unique doit être
nécessaire.

\struck{Et} En effet nous cherchons l'essence ou la nécessité de chaque
chose ; et ces deux \struck{représentent pour} expressions sont
équivalentes ; car\note{« expressions » est écrit au-dessus de la partie
biffée ; le « E » de « En » surcharge « Et ». La phrase se poursuit à la
vue suivante.}

\page{56}
lorsque nous connoissons l'essence, nous voyons que l'être auquel elle
appartient ne sauroit ni n'être pas ni être différent de ce qu'il est.
Notre esprit, satisfait, appuyé sur la nécessité, jouit alors d'une
parfaite quiétude. L'attrait des sciences exactes n'a pas d'autre cause.
Les sujets qu'elles embrassent, sont connus dans leur essence ; leur
existence est tellement nécessaire qu'on ne \struck{peut} sauroit même
concevoir qu'ils puissent ne pas exister. L'esprit se plaît à les
considérer parce qu'il entre ainsi dans l'intime possession de l'être
nécessaire ou de la vérité pure.

Partout ailleurs nous ne voyons plus que des êtres dépendans, des vérités
partielles. Nous cherchons \struck{la cause} l'origine de tels êtres, la
vérité nécessaire dont \struck{elles ci} émanent ces vérités
partielles.\note{« l' » surcharge un « d' » ; le mot biffé après « dont »
est lu avec doute.}

À l'égard des objets de ce genre, nous n'éprouvons aucune répugnance à
admettre qu'ils peuvent ne pas exister : ou, ce qui est une idée
semblable, nous accordons aisément qu'ils pouvoient être différents de ce
qu'ils sont en effet.

Il est évident que cette disposition de notre esprit tient uniquement à
l'ignorance où nous sommes touchant un tel ordre de choses.

Aussi voyons nous, \struck{à mesure} que les progrès des sciences, en
établissant à nos yeux des liaisons entre des faits que nous avions crus
isolés, nous forcent, lorsque les uns sont constatés, à regarder les
autres comme nécessaires. C'est \struck{que nous reconnoissons alors}
qu'alors, nous \struck{regardons} considérons ces faits comme des parties
différentes d'une même existence ; tandis qu'auparavant, nous pensions
qu'ils appartenoient\note{La phrase se poursuit à la vue suivante.}

\page{57}\folio{28r}
à des unités différentes.

Lorsqu'il s'agit de faits éventuels, l'analyse nous sert à calculer dans
un cas donné, la probabilité que tel fait arrive plutôt que tel autre.
Notre réponse à la question de \struck{relative à} la possibilité du
fait, quelque soit la nature de ce fait, est empirique. Sans se mettre en
peine des circonstances qui peuvent en amener la réalisation, le géomètre
répond : il y a une cause pour que tel événement arrive quelque fois ;
la probabilité que cette cause l'amènera \struck{en effet} l'événement
est exprimée par \struck{la} telle fraction \struck{que \ill{}}.\note{Les
deux derniers mots biffés de la ligne ne sont pas lus.}

L'utilité d'une telle réponse est incontestable ; mais elle atteste notre
ignorance. Par exemple, si on demandoit quelle est la probabilité qu'une
certaine machine vienne à cesser à \struck{tel mo} un instant déterminé ;
on auroit sans doute grand intérêt à \struck{connoître en effet} savoir
cette probabilité, mais il est clair que, si on connoissoit parfaitement
la force employée, les frottemens et les résistances, on
\struck{pourroit} sauroit que l'événement est inévitable ou qu'il est
impossible ; et cette impossibilité seroit évidente pour l'instant qui
précède immédiatement celui \struck{pour} lequel la réalisation du même
événement seroit aussi évidente.\note{« lequel » est écrit sur « pour »
biffé et surchargé.}

Dans des événemens d'une nature plus compliquée, nous ne sommes même pas
en état de dire quelles sont les connoissances qui nous manquent pour
acquérir la certitude. Mais parce que nous ignorons quelles sont les
circonstances déterminantes, devons nous penser qu'elles sont
arbitraires, sans liaisons, sans ordre, enfin qu'elles manquent à toutes
les conditions qui \struck{nous} le\note{La phrase se poursuit à la vue
suivante.}

\page{58}
\struck{\ill{}} rencontrent dans toutes les réalités qui sont à notre
connoissance ?\note{Le mot biffé, sous « rencontrent », est effacé d'une
traînée d'encre.}

Conclusons donc que la distinction entre les faits contingens et
nécessaires est, quant au fond, celle qui se trouve entre les faits dont
on ignore et ceux dont on connoît la nature.\note{Un crochet ouvrant est
tracé devant ce paragraphe.}

L'univers, ce fait unique dont l'existence tourmente depuis si longtems
l'esprit des philosophes, s'il étoit mieux connu, paroîtroit nécessaire.
On sait que cette opinion a été soutenue.

Des distinctions entre l'intelligence et la matière, distinctions dont
nous avons signalé l'origine, ont fait \struck{regarder} repousser la
nécessité jusqu'à Dieu ; et l'idée de dieu \struck{s'est trouvée} a été
formée sur le modèle de notre intelligence.

On a dit : dieu est nécessaire ; sa volonté est libre ; il a voulu
l'univers. Mais, en disant : sa volonté est libre, on a rompu la chaîne ;
car, s'il a pû ne pas vouloir l'univers, l'univers n'émane plus de lui
comme les vérités secondaires émanent de l'unité nécessaire dont elles
font partie. Il est clair que le sentiment de liberté qui accompagne les
déterminations de notre volonté a été le modèle qu'on a suivi ; et
pourtant ce sentiment lui même ne peut nous empêcher de reconnoître que
notre volonté est souvent entraînée par les lois imprescriptibles de la
nécessité. Il est vrai que nous délibérions très réellement ; mais nous
nous décidons. Semblables à la balance dont les deux plateaux sont
chargés, nous \struck{oscillons} oscillons ; mais le poids le plus fort
détermine la situation où le système \struck{est} demeure en
repos.\note{« nous délibérions » est écrit ainsi. « demeure » est ajouté
au-dessus de « est » biffé.}

\page{59}\folio{29r}
Il est naturel que la délibération nous donne le sentiment de notre
liberté, et nous distraie même de la prévision d'une détermination, qui
bien que nécessaire, nous semble avoir été sur le point d'être changée en
une détermination contraire. \struck{mais} Aussi voyons que une personne
qui connoît à la fois la position et le caractère d'une autre personne
prévoit avec certitude le parti que \struck{elle} prendra \struck{tandis
que la première} celle-ci, qui, étonnée de cette espèce de prédiction
assure, avec vérité, qu'il s'en est peu fallu qu'elle n'ait agi d'une
façon différente.\note{« Aussi voyons » est écrit ainsi. « celle-ci,
qui, » est ajouté au-dessus de « que la première », dont « que » ne
paraît pas biffé.}

Plus on réfléchit, plus aussi on reconnoît que la nécessité gouverne le
monde. À chaque progrès nouveau des sciences, ce qui passoit pour
contingent est reconnu comme étant nécessaire. Il s'établit à nos yeux
des liaisons multipliées entre des branches qu'on avoit crues séparées ;
on observe des lois, là où l'on n'avoit encore vu que des faits
accidentels. Nous approchons de plus en plus de l'unité d'être, qui fut
le rêve de l'antiquité, et qui trouve son modèle dans le sentiment de
notre propre existence.

\uncertain{Tâchons} enfin de fixer notre opinion à l'égard de ce modèle
du vrai, de ce type de l'être qui a souvent égaré la raison humaine ; et
qui, dans nos tems modernes, sert à la guider d'une manière si heureuse,
que ses progrès, \struck{de la} d'abord concentrés entre un petit nombre
d'hommes adonnés à l'étude, se répand aujourd'hui dans toutes les classes
de la société ; \struck{et} éclaire à la fois, les sciences morales
politiques, la physique, les arts chimiques et mécaniques et peut fournir
encore aux lettres et aux beaux arts des lumières nouvelles, des
inspirations qu'ils n'ont pas encore rencontrées.\note{Le premier mot du
paragraphe est récrit sur un autre, et la surcharge le rend douteux.
« ses progrès … se répand » est écrit ainsi.}

\page{60}
\note{En tête du feuillet, au crayon et d'une autre main, une ligne très
pâle dont on lit seulement « Ch 4 » et, plus loin, « type »,
\ill{}.}
L'homme, \struck{s'il n'avoit} peut-il pas d'autre sujet d'étendue que
lui même, connoîtroit l'étendue ; je ne pense pas qu'il puisse
sérieusement douter de cette propriété de la matière mais ce qu'il
connoîtroit surtout avec la dernière évidence, c'est sa propre
existence.\note{« peut-il » est écrit au-dessus de « s'il n'avoit »
biffé ; la phrase est laissée telle quelle.}

\struck{Dans les} Au milieu des divers systèmes où s'est égaré l'esprit
humain, il a essayé du scepticisme. Il a pû soutenir que tout ce qui
étoit au dehors de l'existence de l'homme étoit pure apparence ; mais le
fameux argument, « je pense, donc je suis » a ramené la raison vers la
réalité de son être.\note{Un crochet ouvrant est tracé devant ce
paragraphe.}

Le sentiment de l'être est celui de la vérité. Il est inséparable de
notre existence, il précède toute autre idée. Le bon, le beau dérivent du
vrai ; mais leur connoissance exige le secours des comparaisons. Suivant
qu'on s'est trouvé plus ou moins frappé de l'une ou de l'autre des parties
de cette \struck{assertion} proposition, on a été porté, par l'esprit de
systèmes, à soutenir, ou que nos idées sont innées, ou qu'elles viennent
de nos sensations. L'une et l'autre opinions sont \struck{\ill{}} vraies
dans les limites que nous avons posées.\note{« ou de l'autre » est ajouté
dans la marge de gauche. « opinions » est ajouté au-dessus de la ligne,
avec un renvoi ; « vraies » est écrit au-dessus d'un mot biffé et
barbouillé.} Le type du vrai, nous l'apportons en naissant, notre être
dont la réalité est notre plus intime connoissance est inséparable de ce
modèle inné. En ce sens, l'homme est l'abrégé de l'univers ; car l'être
ou la vérité, partout où ils se trouvent, remplissent certaines
conditions, que l'attention découvrira nécessairement \struck{\ill{}} dans
tous les objets réels dont elle sera occupée.

Mais cette ressemblance abstraite est fort éloignée de celle qu'on a
cherchée. Elle peut servir à expliquer cependant la cause d'une
erreur\note{« servir » est couvert d'une tache d'encre. La page s'arrête
ici, au milieu de la phrase.}

\end{document}
